\documentclass{report}
% Ćirilica: kodiranje fonta i jezik (srpski, ćirilica)
\usepackage[T2A]{fontenc}
\usepackage[utf8]{inputenc}
\usepackage[serbianc]{babel}

% Veličina strane i margine
\usepackage[a4paper,top=2cm,bottom=2cm,left=3cm,right=3cm,marginparwidth=1.75cm]{geometry}

% Korisni paketi
\usepackage{amsmath}
\usepackage{graphicx}
\usepackage{float}
\usepackage{booktabs}
\usepackage{enumitem}
\usepackage{tcolorbox}
\tcbuselibrary{breakable}
\usepackage[colorlinks=true, allcolors=blue]{hyperref}

% Нумерација као у уџбенику: 3, 3.1, 3.1.1, 3.1.1.1
\setcounter{secnumdepth}{4}
\setcounter{tocdepth}{3}
\setcounter{section}{2}

% Оквир за кључне податке
\newtcolorbox{kljucno}{colback=blue!4, colframe=blue!45!black,
    boxrule=0.5pt, arc=1.5pt, left=8pt, right=8pt, top=4pt, bottom=4pt}

% Слика са провером: ако фајл још не постоји, црта се празан оквир са именом фајла
%   #1 -- ширина (део ширине текста), #2 -- име фајла у фолдеру slike, #3 -- натпис
\newcommand{\slika}[3]{%
\begin{figure}[H]
    \centering
    \IfFileExists{slike/#2}%
        {\includegraphics[width=#1\textwidth]{slike/#2}}%
        {\fbox{\parbox[c][3cm][c]{#1\textwidth}{\centering\small слика: \texttt{slike/#2}}}}
    \caption{#3}
\end{figure}}

\title{Аутоелектронски системи 2 -- Осветљење и сигнализација}
\author{Саобраћајно техничка школа, Београд, Земун}
\date{}

\begin{document}
\maketitle
\tableofcontents
\newpage


%==================================================================
\subsection{Прање и брисање ветробрана и фарова}
%==================================================================

\subsubsection{Прање и брисање ветробрана и задњег стакла}

Чиста и провидна стакла предуслов су за добру видљивост, а тиме и за
безбедну вожњу, нарочито ноћу. Зато уређаји за брисање морају да обезбеде:
\begin{itemize}[noitemsep]
    \item да брисана површина ветробрана буде што већа и што чистија,
    \item квалитет брисања који спречава расипање светлости и заслепљивање,
    \item да брисачи раде тихо и поуздано, без застоја, и да имају
          интервални рад или сензор за кишу.
\end{itemize}

Такозваним „пери-бриши” уређајем стакла се перу и бришу, напред и назад, у
свим временским условима. Важан је и квалитет гумице брисача: данас се
израђују у двослојној техници, са меком гумом и оштрим ивицама.

\slika{0.6}{sl3-59.png}{Комплетан „пери-бриши” уређај за савремено путничко возило: 1. „пери-бриши” уређај за фарове и предњи ветробран, 2. предња посуда за воду са пумпом, 3. брисаљке за предњи ветробран, 4. брисачи задњег стакла, 5. задње брисаљке, 6. „пери-бриши” уређај за задњи ветробран, 7. задња посуда за воду са пумпом, 8. уређај за прање фарова високог притиска, 9. посуда за воду, 10. пумпа високог притиска}

\slika{0.6}{sl3-60.png}{Електрична шема повезивања комплетног „пери-бриши” уређаја: M4 -- мотор брисача, M5 -- мотор пумпе за прање ветробрана, M6 -- мотор брисача фарова, M7 -- мотор пумпе за прање фарова, K8 -- релеј за интервални рад брисача, S27 -- прекидач брисача ветробрана, S28 -- прекидач за прање и брисање фарова, S29 -- прекидач за прање ветробрана}

Детаљан опис механизма брисача и уређаја за интервално укључење брисача дат
је у уџбенику за трећи разред (Електрични уређаји на моторним возилима).

\subsubsection{Уређај за прање и брисање предњих фарова}

Овај уређај чисти запрљана стакла фарова, чиме се побољшава осветљење пута и
смањује заслепљивање возила из супротног смера. У употреби су два система:
„пери-бриши” уређаји са брисачима и прскалице високог притиска.

„Пери-бриши” уређај са брисачима и прскалицама ради слично као код
ветробрана, али има недостатке: зими се брисачи тешко одлепе од смрзнутог
стакла, а код нових фарова са пластичним стаклима (поглавље 3.5.2) брисачи
могу да изгребу стакло. Зато се данас фарови чисте само прскалицама високог
притиска.

\paragraph*{Прскалице високог притиска}
(слика 3.61) перу стакло јаким млазом воде. Квалитет прања зависи од
растојања између млазнице и стакла, угла под којим млаз пада, величине и
брзине капљица и количине утрошене воде. Млазнице се постављају на браник
или на телескопски носач, који се при прању извлачи из браника.

Према европским прописима од 1996. године, сва нова возила са гасним
сијалицама морају имати уређај за прање фарова. Класа 25 значи да количина
воде у посуди довољна за 25 циклуса прања, а класа 50 за 50 циклуса. Прање
мора да смањи губитак светлосног флукса због прљавштине са 70\,\% на највише
20\,\%. Уређај мора да ради на температурама од $-35$ до $+80$\,$^\circ$C, па се
зими користе додаци против смрзавања, а понекад и загревање.

\slika{0.4}{sl3-61.png}{Главни делови уређаја за прање фарова прскалицама високог притиска: 1. посуда за воду, 2. пумпа, 3. неповратни вентил, 4. Т-рачва за разводњавање воде за оба фара, 5. држач млазнице, 6. црева}

\subsubsection{Сензори за кишне капи и прљавштину}

\paragraph*{Сензор за кишне капи}
(слика 3.62) је оптички уређај са предајником и пријемником светлости,
постављен на ветробрану са унутрашње стране. Светлосни зраци се под
одређеним углом усмеравају ка спољној страни ветробрана. Када је стакло суво,
зраци се на спољној површини потпуно одбијају (тотална рефлексија) и враћају у
пријемник. Када на стакло падну капи воде, услови рефлексије се мењају, део
зрака излази напоље, па је сигнал у пријемнику слабији. Сензор тада
аутоматски укључује брисаче.

\slika{0.45}{sl3-62.png}{Сензор за кишне капи на ветробрану: 1. капљице кише, 2. стакло ветробрана, 3. огледало, 4. бленда, 5. сензор одбијене светлости, 6. извор светлости, 7. огледало}

\paragraph*{Сензор за прљавштину}
(слика 3.63) се поставља у фар, са унутрашње стране стакла, поред светлосног
снопа. Састоји се од извора светла (најчешће LED-диоде, у инфрацрвеном
подручју) и пријемника (фото-транзистора). Када је стакло чисто, зраци
пролазе напоље. Честице прљавштине на стаклу одбијају део зрака назад у
пријемник, сразмерно степену запрљаности. Када запрљаност пређе одређену
границу, сензор, ако су светла укључена, аутоматски укључује прање фарова.

\slika{0.45}{sl3-63.png}{Сензор за прљавштину: 1. честице прљавштине на стаклу, 2. стакло фара, 3. пријемник, 4. извор светла, 5. кућиште сензора}

%==================================================================
\subsection{Звучни сигнални уређаји}
%==================================================================

По прописима свако возило мора имати бар један уређај за звучне сигнале
(сирену), којом у одређеним ситуацијама упозорава остале учеснике на
опасност. Звук сирене мора се разликовати од буке возила и средине у којој
се возило налази, па је важан добар избор тона. Највећи домет (300 до
400~m) имају звучни сигнали учестаности мање од 1000~Hz, јер их ваздух
најмање упија, а сигнали испод 500~Hz продиру и кроз чврсте објекте, па се
боље чују у затвореним возилима.

Ако возило има више сирена, по прописима је забрањена њихова истовремена
употреба. Јачина звука, мерена 7~m испред возила на висини од 500 до
1500~mm, не сме бити већа од 105~dB. Посебна возила (хитна помоћ, полиција,
ватрогасци) имају звучне сигнале који мењају висину звука или раде са
прекидима, јер најбоље привлаче пажњу.

Постоје електричне сирене са мембраном, електропнеуматске трубе и моторне
пнеуматске трубе.

\subsubsection{Електрична сирена са мембраном}

Ради слично као електрично звоно. Електромагнет се напаја преко малог
прекидача у самој сирени, чији су контакти механички спојени са котвом
магнета, тако да магнет сам себе искључује (слика 3.64). Када потече струја,
магнет привуче котву, котва отвори контакте и прекине струју. Опруга тада
врати котву, контакти се поново затворе и процес се понавља. Кондензатор
везан паралелно са контактима смањује варничење и продужава им век.

Котва је спојена са еластичном мембраном и звучном плочом, па цео склоп
трепери великом учестаношћу. Мембрана даје основни тон, учестаности од 250
до 400~Hz, а звучна плоча више тонове (од око 1800 до 2000~Hz). Јачина и боја
звука зависе од материјала, облика и димензија мембране и плоче, а
подешавају се променом размака контаката. Сирена се поставља на место које
неће треперити заједно са њом, јер би се тиме смањила јачина звука.

\slika{0.4}{sl3-64.png}{Принцип рада електричне сирене са мембраном (шематски приказ): 1. плочасти одушак, 2. звучна плоча, 3. еластична мембрана, 4. котва магнета, 5. електромагнет, 6. контакти прекидача, 7. кондензатор, 8. еластични држач}

\subsubsection{Електропнеуматска труба}

Даје чист и пријатан звук, сличан трубама. Овде мембрана не производи звук
директно, већ побуђује на треперење ваздушни стуб у спирално савијеној цеви
(слика 3.65). Мембрана нема звучну плочу. Висина звука зависи од учестаности
треперења мембране, која зависи од резонантне учестаности ваздушног стуба.
Цев је дуга, савијена у спиралу ради уштеде простора, и постепено се шири
као левак, ради бољег преноса звука.

Звук трубе садржи цео спектар виших тонова (од 600 до 1000~Hz). Звук је
уједначенији и мекши него код сирене са мембраном, али слабије продире кроз
буку. Зато се трубе обично постављају у паровима различите висине, па се
често зову „двотонске сирене”.

\slika{0.45}{sl3-65.png}{Шематски приказ електропнеуматске трубе}

\subsubsection{Моторна пнеуматска труба}

Електромотор покреће мали вентилатор, који потискује ваздух у комору са
еластичном мембраном. Мембрана затвара улаз у трубу. Када притисак порасте,
мембрана се помери, ваздух струји кроз трубу, притисак опадне и мембрана се
врати, па настаје треперење. Боја звука зависи од облика и величине трубе.
Недостатак им је што не реагују одмах, јер вентилатор мора прво да створи
потребан притисак.

\subsubsection{Електрично напајање сирена}

За рад сирене веома је важно да напајање буде квалитетно и да пад напона у
водовима буде што мањи. Због велике снаге сирене, водови морају бити
дебели и што краћи. Зато није погодно да цела струја иде од акумулатора,
преко осигурача и тастера на волану, до сирене. Најчешће се сирена напаја
\textbf{преко релеја} (слика 3.66).

Сирена се дебелим водовима везује директно, преко релеја и осигурача, за
извор сталне струје (30), чиме се водови знатно скраћују. Тастер на волану
укључује само намотај релеја, кроз који тече мала струја, па водови до
тастера могу бити танки. Сирена се, дакле, не укључује директно преко
тастера, већ индиректно, преко контаката релеја, који имају мали прелазни
отпор и мали пад напона.

\slika{0.35}{sl3-66.png}{Електрично напајање сирене преко релеја: 1. тастер сирене, 2. осигурач, 3. релеј, 4. сирена, 30. стална струја од акумулатора, 85 и 87. контакти релеја}

\subsubsection{Алармне сирене}

Хитна помоћ, ватрогасци и полиција у хитним интервенцијама имају предност у
саобраћају, па су опремљени алармним сиренама и ротирајућим упозоравајућим
светлима различитих боја. Уређај се састоји од јаких звучних сирена са дубоким
и високим тоновима и ротирајућих светала. На слици 3.67 приказана је шема
везе алармних сирена и ротирајућих светала преко командног уређаја. Командни
уређај у одређеним временским размацима наизменично укључује сирене, а
ротирајућа светла укључује истовремено.

\slika{0.5}{sl3-67.png}{Шема везе алармних сирена и ротирајућих светала преко командног уређаја: 1. алармне сирене, 2. нормална сирена, 3. командни уређај, 4. прекидач аларма, 5. контролна лампица, 6. ротирајућа светла, 7. осигурач, 8. тастер сирене}

Командни уређај (слика 3.68) састоји се од астабилног мултивибратора, који
одређује учестаност укључивања сирена, и моностабилног мултивибратора, који
одређује трајање импулса (на пример $4 \times 0{,}75~\text{s} = 3$~s). Ротирајућа
светла раде преко контролних степена, слично као упозоравајући мигавци.

\slika{0.55}{sl3-68.png}{Електронски командни уређај BOSCH за акустично-оптичке аларме: 1. алармне сирене, 2. ротирајућа светла, 3. контролне лампице, 4. астабилни мултивибратор, 5. моностабилни мултивибратор, 6. контролни степен}

%==================================================================
\subsection{Контролни инструменти}
%==================================================================

Контролни инструменти дају возачу информације о стању возила и раду
најважнијих агрегата током вожње. Тиме се повећава безбедност, вожња је
економичнија, а кварови се откривају на време. Безбедност, економичност и
еколошки захтеви су све важнији, па се, уз развој електронике, систем
контроле стално усавршава и обухвата све више величина.

\subsubsection{Брзиномер и километар-сат}

Брзина кретања, коју показује брзиномер, једна је од најважнијих информација
за возача. Без ње је тешко прилагодити се условима вожње и прописима. Брзина
којом се може безбедно возити зависи од:
\begin{itemize}[noitemsep]
    \item техничких карактеристика, димензија и оптерећења возила,
    \item временских услова и видљивости на путу,
    \item врсте пута, стања коловоза и квалитета пнеуматика,
    \item густине саобраћаја и ограничења брзине,
    \item доба дана (дању или ноћу) и искуства возача.
\end{itemize}

\paragraph{Опис и конструкција механичког брзиномера}

На почетку развоја брзина се мерила механичким инструментима са
центрифугалним теговима. Касније су уведени знатно бољи
електромеханички брзиномери са индукционим мерним системом, који су и
данас у масовној употреби. У новије време све више се користе потпуно
електрични или електронски брзиномери.

\textbf{Индукциони брзиномер} (слика 3.69). На обртну осовину (8) причвршћен
је вишеполни стални магнет у облику диска (1). Осовину окреће механичка
сајла, повезана са мењачем, па се магнет окреће брзином сразмерном брзини
возила. Изнад магнета је, без додира са њим, мерно звоно (7) од немагнетног
али електрично добро проводног материјала (најчешће алуминијума). Звоно је на
осовини (6), на којој је и казаљка (4).

Док се магнет окреће, његове линије поља (2) секу звоно и у њему индукују
вртложне (Фукоове) струје. Њихово магнетно поље делује са пољем магнета, па
настаје обртни момент који повлачи звоно у смеру обртања магнета. Том
моменту супротставља се спирална опруга (3), па се звоно окрене само толико
док се моменти не изједначе. Угао заокретања звона, а тиме и положај
казаљке на скали (5), сразмеран је брзини возила.

\slika{0.4}{sl3-69.png}{Принцип рада индукционог брзиномера: 1. стални магнет у облику диска, 2. линије магнетног поља, 3. спирална опруга, 4. казаљка, 5. скала, 6. осовина мерног звона, 7. мерно звоно, 8. обртна осовина}

Недостатак овог уређаја је велика зависност од температуре: са порастом
температуре слаби магнет и расте отпор звона, па грешка у екстремним
условима може бити и већа од 15\,\%. Зато се користи специјални
термомагнетни шант, чија магнетна својства зависе од температуре, па се
грешка смањује на око 1,5\,\% на 10\,$^\circ$C.

\paragraph{Опис и конструкција електричног брзиномера}

\textbf{Показни инструмент.} Код електричних брзиномера користи се
инструмент са покретним калемом (слика 3.70). Има широку скалу ($270^\circ$)
и отпоран је на вибрације. Ради као и обични инструменти са покретним калемом:
стални магнет (5), између наставака за полове, ствара хомогено магнетно поље
(8). У пољу је покретни калем (7), повезан са казаљком (1) и спиралним
опругама (3). Када кроз калем тече струја, на проводнике делује
електромагнетна сила која окреће калем, а опруге му се супротстављају.
Отклон казаљке сразмеран је струји, а пошто је поље хомогено, скала је
линеарна.

\slika{0.45}{sl3-70.png}{Принцип рада инструмента са покретним калемом: 1. казаљка, 2. осовина, 3. спирална опруга, 4. наставак у облику цилиндра (северни пол), 5. стални магнет, 6. полни наставак у облику цилиндра (јужни пол), 7. покретни калем, 8. линије хомогеног магнетног поља}

\textbf{Давач електричног брзиномера.} Уместо механичке сајле до
инструмента, мали генератор (тахо-динамо) на мењачу даје напон сразмеран
броју обртаја. Најчешће је то генератор наизменичне струје са обртним
сталним магнетом и непомичним намотајима статора, чија су учестаност и
напон сразмерни брзини возила. Наизменични напон се исправља диодама у
мосној (Грецовој) вези у самом инструменту (слика 3.71). Струја се ограничава
предотпором ($R_e$), а нелинеарност диода утиче на скалу само на почетку.
Температура утиче на карактеристике диода, али се употребом већег
предотпора тај утицај смањује.

\slika{0.4}{sl3-71.png}{Напајање показног инструмента електричног брзиномера преко диода у мосној вези и превођење наизменичне струје у једносмерну}

\paragraph{Механички километар-сат}

Осовина брзиномера може да покреће и километар-сат: бројчаник са шест
међусобно повезаних добоша са бројевима од 0 до 9 (слика 3.72). Први добош
(5) покреће се преко вишестепеног пужастог преноса (2, 3 и 4).

\slika{0.45}{sl3-72.png}{Шематски приказ механичког покретања бројчаника километар-сата: 1. погонска осовина, 2, 3 и 4. пужасти преноси, 5. почетни добош, 6. покретач сајле са великим преносом, 7. главни пренос}

\textbf{Бројчаник.} Конструкција омогућава да се добош окрене за део круга
тек када претходни направи цео круг (слика 3.73). Сваки добош има са једне
стране 20 зуба, а са друге само два. Између два суседна добоша налази се мали
преносни зупчаник са шест зуба. Док први добош прави цео круг, његова два
зуба само једном захвате преносни зупчаник и окрену га за трећину круга (два
зуба). Тај зупчаник окрене следећи добош за два од његових 20 зуба, тј. за
1/10 круга, што је тачно један број. Исто важи и за све остале добоше. Тако
се сваки добош помери за један број тек када претходни направи пун круг.
После 999\,999,9~km бројчаник се враћа на нулу. Посебна конструкција
спречава да се добоши сами помере због вибрација током вожње.

\slika{0.45}{sl3-73.png}{Принцип рада бројчаника са спољашњим озубљењем}

\paragraph{Електрични километар-сат}

Код возила са електричним брзиномером бројчаник се не може покретати
механички, па се користи мали корачни мотор или електромагнетни релеј.
Импулси се добијају од малог генератора или од прекидача на погонској
осовини. На слици 3.74 приказан је прекидач фирме VDO са тренутним
прекидањем: опруга обезбеђује нагло отварање и затварање контаката.

\slika{0.3}{sl3-74.png}{Прекидач фирме VDO са тренутним прекидањем контаката}

\textbf{Бројчаник са релејем.} При сваком импулсу релеј привуче котву и
затегне опругу. Када импулс престане, опруга врати котву, која преко
механизма окрене први добош бројчаника за одређени угао. Однос преноса и
учестаност импулса подешени су тако да први добош направи цео круг на
сваких 1000~m.

\textbf{Бројчаник са корачним мотором.} Покретање и подешавање односа преноса
знатно су једноставнији са малим корачним (степ-)мотором, који се при сваком
импулсу окрене за тачно одређени угао (корак). Мотори за километар-сат се при
сваком импулсу окрену за два корака од по $60^\circ$, тј. за $120^\circ$ (1/3
круга).

Ротор корачног мотора је стални магнет, а статор има три намотаја (слика
3.75). Средњи намотај је стално под напоном и ствара магнетно поље које држи
ротор у мирујућем положају. Два спољна намотаја добијају импулсе, који мењају
поларитет наставака статора, па се ротор окреће за један корак. Облик
наставака статора и полова ротора (слика 3.76) одређује смер обртања.

\slika{0.35}{sl3-75.png}{Шема везе статора корачног мотора фирме VDO: 1. довод импулса, 2. намотаји статора}

\slika{0.4}{sl3-76.png}{Полови сталног магнета ротора (горе) и посебно обликовани наставци полова статора: 1. распоред полова ротора, 2. главни полови статора, 3. помоћни полови статора}

\paragraph{Електронски брзиномер и километар-сат}

Развој електронике постепено потискује механичке и електромеханичке
брзиномере. Савремена возила имају електронске даваче и електронска кола за
обраду сигнала, који су поузданији, трајнији и мањи. Сигнал брзине користе и
други уређаји: путни рачунар, регулатор брзине (темпомат) и слично. На
слици 3.77 дат је блок-дијаграм електронског брзиномера и километар-сата.

Сигнали давача (Reed-контакта или електронског давача) прво пролазе кроз
пропусни филтар, који уклања сметње, а затим их Schmitt-trigger обликује у
правоугаоне импулсе исте амплитуде. Њима се окида моностабилни
мултивибратор, који даје импулсе сталне ширине. Средња вредност тих импулса
(добијена у сабирном колу) сразмерна је учестаности, тј. брзини, и покреће
инструмент са покретним калемом. Бројчаник покреће мали корачни мотор, коме
се учестаност импулса претходно смањи (на пример 32:1). Дигитални излази
служе за сигнализацију прекорачења брзине или за аутоматско ограничење
највеће брзине.

\slika{0.65}{sl3-77.png}{Блок-шема електронског брзиномера и километар-сата са електричним и механичким елементима: 1. Reed-контактни давач, 2. електронски давач, 3. осцилаторни окидач, 4. пропусни филтар, 5. Schmitt-trigger, 6. моностабилни мултивибратор, 7. сабирно коло, 8. сигнализација прекорачења брзине}

\subsubsection{Обртомери}

Код мотора са унутрашњим сагоревањем веома је важно пратити број обртаја.
Радом у подручју оптималног обртног момента смањује се потрошња горива,
продужава век мотора и правилније се бира степен преноса. Обртомер је
нарочито важан код спортских и теретних возила, где возач не може да
процени број обртаја по звуку мотора.

\paragraph{Претварачи сигнала -- давачи}

Као давачи се користе мали генератори, индуктивни претварачи и систем
паљења.

\textbf{Електрични генератор} ради као тахо-динамо код брзиномера.

\textbf{Индуктивни давач} (слика 3.78). Импулси се индукују у малом намотају
(1) са језгром од меког гвожђа (2), који се налази у пољу сталног магнета (3),
поред зупчастог точка од меког гвожђа (4). Када поред језгра пролази зуб,
мења се магнетни флукс, па се у намотају индукује напон. Учестаност импулса
сразмерна је броју обртаја и броју зуба. Амплитуда импулса је мала, па се
појачава и обликује електронским колом, обично у самом инструменту.

\slika{0.35}{sl3-78.png}{Индуктивни давач сигнала за обртомер (VDO): 1. намотај, 2. језгро од меког гвожђа, 3. стални магнет, 4. зупчасти точак од меког гвожђа}

\textbf{Систем паљења као давач.} Код већине електронских обртомера као
улазни сигнал користе се импулси из примарног кола система паљења, чија је
учестаност сразмерна броју обртаја. Уређаји за дијагностику користе и
импулсе из секундарног кола.

\paragraph{Транзисторски обртомери}

Електронски обртомери као улазне сигнале користе импулсе са прекидача
паљења. Ти импулси, због самоиндукције, садрже осцилације и нису
правилног облика, па се морају обликовати. Најчешће се то ради
моностабилним мултивибратором, који за сваки улазни импулс даје импулс сталне
ширине и висине. Средња вредност струје тих импулса сразмерна је броју
обртаја и покреће показни инструмент.

\textbf{Обртомер TH 193} (слика 3.79). Улазни сигнали (тачка 1) пролазе кроз
филтре ($R_1$--$C_1$ и $R_2$--$C_2$), а диода $D_1$ и кондензатор $C_3$
одсецају осцилације, па до базе транзистора стижу само врхови импулса
(тачке 2 и 3). Транзистори $T_1$ и $T_2$ чине моностабилни мултивибратор,
који за сваки импулс даје импулс сталног трајања, одређеног RC члановима.
Показни инструмент добија средњу вредност струје, сразмерну броју обртаја.

\slika{0.5}{sl3-79.png}{Шема везе транзисторског обртомера за путничка возила TH 193, са изгледом облика сигнала после обраде у тачкама 1, 2 и 3}

\textbf{Обртомер MS-7} (слика 3.80) ради на истом принципу. NTC отпорник
компензује утицај температуре, ценер-диода стабилизује напон напајања, а
промењивим отпорником се инструмент подешава (калибрише).

\slika{0.6}{sl3-80.png}{Шема везе електронског обртомера MS-7 за путничка возила}

Показни инструмент ради на принципу покретног калема са сталним магнетом и
има скалу од $270^\circ$. Струја пуног скретања инструмента је око 10~mA.

Шема обртомера MS-7 донекле се разликује од претходне. Кондензатор $C_1$
пропушта само наизменичну компоненту сигнала, а са отпорницима ствара степен
пасивног диференцијатора, који од улазних осцилација добија кратке импулсе.
Диода одсеца негативне импулсе, а само позитивни доспевају до базе
транзистора. Транзистори $T_1$ и $T_2$ чине моностабилни мултивибратор, који
за сваки улазни импулс даје правоугаони импулс сталне ширине и амплитуде.

Посебност ове везе је у томе што показни инструмент није у колу једног од
транзистора мултивибратора, већ у колекторском колу посебног транзистора
($T_4$). Тако се постиже боља напонска и температурна стабилизација. Отпорник
($R_{11}$) са негативним температурним коефицијентом (NTC) компензује утицај
температуре, а диода ($D_4$) између базе и емитера транзистора $T_4$
компензује температурну зависност напона база--емитер. Ценер-диода ($D_5$),
заједно са транзистором ($T_3$), стабилизује напон напајања, а
потенциометром ($R_{13}$) подешава се отклон инструмента. На доњем делу
слике 3.80 приказани су облици сигнала у појединим тачкама везе, који се
могу контролисати осцилоскопом.

\textbf{Показни инструмент.} На слици 3.81 приказана је конструкција
инструмента са покретним калемом и скалом од $270^\circ$. Калем (6) обрће се у
магнетном пољу сталног магнета (3), између наставака полова (2 и 10), од којих
је (10) у облику прстена у коме се налази калем. Осовина (4) лежи у
лежиштима (5 и 9). Казаљка (8) на скали (7) показује отклон калема, који је
сразмеран средњој вредности струјних импулса.

\slika{0.45}{sl3-81.png}{Шематски приказ конструкције инструмента са крећним калемом: 1. носећа конструкција, 2. наставци полова, 3. стални магнет, 4. осовина инструмента, 5 и 9. лежишта, 6. крећни калем, 7. скала, 8. казаљка, 10. гвоздено језгро}

%------------------------------------------------------------------
\subsubsection{Тахографи}
%------------------------------------------------------------------

Тахографи су регистрирајући инструменти који писаљкама бележе брзину возила,
пређени пут, време вожње и стајања и друге важне информације. Уграђују се у
камионе и аутобусе, заједно са брзиномером и километар-сатом.

\textbf{Папир за регистровање} је кружни дијаграм са концентричним линијама,
на коме писаљке инструмента остављају траг. Пошто се подаци бележе у
зависности од времена, папир покреће сатни механизам, обично за цео круг
за 24 сата (слика 3.82). По спољашњем ободу дијаграма је сатна скала, са
поделом на 5 минута. Брзина се бележи радијалним скретањем писаљке, по скали
са поделом од 20~km/h. Пређени пут бележи се у унутрашњем кругу, у облику
троуглова: сваки „троугао” одговара путу од 10~km. Праве линије означавају
време стајања, а дебела линија време вожње. У средини дијаграма возач
ручно уписује основне податке: име, полазно и одредишно место, датум и
стање километар-сата на почетку и крају вожње.

\slika{0.5}{sl3-82.png}{Облик дијаграма на папиру за регистровање у тахографу}

\textbf{Индукциони инструмент.} Код механичких тахографа скоро искључиво се
користе индукциони инструменти, исти као код брзиномера (слика 3.69). Разлика
је само у томе што се користи јачи стални магнет, који даје већи обртни
момент. Мерно звоно преко малог зупчастог сегмента помера писаљку са
сафирним врхом, која бележи податке о брзини на папиру. Писаљку пређеног пута
покреће посебна мала осовина са бројчаника километар-сата, тако да се
окрене за цео круг на сваких 10~km.

\paragraph{Електрични (електронски) тахографи}

Механичка сајла има кратак век због великих статичких и динамичких
оптерећења, па се и код тахографа примењује електрично покретање. Уместо
индукционог мерног система користи се електрични, сличан електричном
брзиномеру: обртање се најпре претвара у електрични сигнал, а у инструменту
се поново претвара у механички момент који покреће казаљку и писаљке.

\textbf{Принцип рада.} Као и код брзиномера, потребна су два претварача:
један за брзину, а други за пређени пут. Давач брзине је мали електрични
генератор чији ротор има 36 полова сталног магнета, а статор такође 36
полова, па даје наизменичну струју од 18 периода по обртају. Најмања
учестаност, од око 15~Hz, постиже се већ при брзини од око 3~km/h. Давач
импулса за пређени пут најчешће је Reed-контакт, који за сваки обртај даје
један импулс, поуздано и при малим брзинама. Давачи се напајају са
акумулатора. Сигнали давача пута доводе се до малог корачног мотора у
инструменту, који покреће бројчаник и механизам писаљке пређеног пута.

Сигнал давача брзине Schmitt-trigger обликује у правоугаоне импулсе, а
моностабилни мултивибратор их уједначава по ширини и амплитуди. Средња
вредност тих импулса сразмерна је брзини и представља електрични еквивалент
момента који треба претворити у механичко скретање. Казаљку и писаљку
брзине покреће мали серво-мотор једносмерне струје, преко појачавача и
управљачког кола, тако да угао скретања буде сразмеран напону. Стварно
скретање региструје мали потенциометар повезан са осовином, чији се сигнал
враћа на улазни диференцијатор. Свако одступање између задатог и стварног
скретања враћа се, као електрични сигнал, преко повратне спреге у
управљачко коло, па се систем окреће док одступање не постане нула. Тако
се повећава тачност и постиже боља статичка и динамичка регулација.
Обртање папира за регистровање покреће сатни механизам са кварцним
осцилатором, ради веће тачности.

\slika{0.55}{sl3-83.png}{Принцип рада електричног тахографа са брзиномером (Kienzle): 1. давач, 2. обликовање сигнала, 3. константни извор, 4. диференцијатор сигнала, 5. појачивач, 6. серво-мотор, 7. повратна спрега, 8. сензор угла померања, 9. управљачко коло, 10. зупчаник, 11. зупчаста летва, 12. писаљка, 13. казаљка брзиномера, 14. скала}

%------------------------------------------------------------------
\subsubsection{Мерачи притиска -- манометри}
%------------------------------------------------------------------

За разлику од брзиномера и километар-сата, који дају податке о кретању
возила, мерачи притиска спадају у инструменте за погонску сигурност возила.
Најчешће мере притисак уља у мотору, потпритисак у усисној цеви (код
SUS-мотора), који даје податке о потрошњи горива, и притисак ваздуха у
кочионом систему, односно у резервоарима ваздуха.

Као и остали инструменти, и ови се састоје од два дела: давача и показног
инструмента. Давачи су претварачи неелектричних величина у електричне и
повезани су са местом мерења. За мерење притиска најчешће се користе давачи
са биметалном траком или потенциометром (слика 3.84).

\slika{0.5}{sl3-84.png}{Шематски приказ давача за мерење притиска: а) са биметалом и б) са потенциометром}

Доњи део давача је од масивног метала и уврће се у простор под притиском, тако
да притисак делује на еластичну мембрану. У горњем делу је електрични део:
биметалне траке са загреваним влакном или потенциометар са променљивим
отпорником.

\textbf{Биметални давач} има горњу биметалну траку са загреваним влакном и
доњу траку, која је у додиру са еластичном мембраном. Контакти трака су
нормално спојени. Када се укључи контакт-кључ, струја тече кроз показни
инструмент, загревано влакно и контакте до масе. Горња трака се загрева,
савија навише, отвара контакте и прекида струју. Када се охлади, контакти се
поново споје. Тако настају струјни импулси.

Ритам импулса зависи од притиска мембране на доњу траку. При притиску нула
контакти се лако и брзо отварају, па су импулси ретки и кратки (слика 3.85,
положај 0). Када притисак порасте, на пример на 2~bar, мембрана јаче потискује
доњу траку, па горња трака мора више да се савије да би отворила контакте.
То захтева дуже загревање, па су импулси дужи (положај 2~bar). При још већем
притиску (5~bar) импулси су још дужи (положај 5~bar).

\slika{0.7}{sl3-85.png}{Принцип рада биметалног мерача притиска са дијаграмом мерене струје у зависности од притиска: $h_1$ и $h_2$ -- висина подизања контакта, $I_0$ -- величина импулса струје, $I_{ef}$ -- ефективна вредност импулса струје, $t_0$ -- време отварања контакта, $t_z$ -- време затварања контакта}

\textbf{Биметални инструмент.} На доњем делу слике 3.85 симболички је
приказано скретање биметалног инструмента у зависности од учестаности
струјних импулса и њихове ефективне вредности. Основни део инструмента је
биметал (2) у облику слова U (слика 3.86). На једном краку налази се загревано
влакно (3), а други, савијени крак повезан је са подешивачем и, преко
назубљеног лука (4), са казаљком (9). Загревано влакно напаја се преко
прикључака (5): један је повезан са биметалним давачем, а други, преко
контакт-кључа, са позитивним полом акумулатора. Скретање казаљке зависи од
ефективне вредности струјних импулса.

\slika{0.45}{sl3-86.png}{Принцип рада и конструкција биметалног инструмента: 1. основна лимена плоча, 2. биметал у облику слова U, 3. влакно које се загрева, 4. назубљени лук подешивача, 5. прикључни контакти за довод струје, 6. лимени подешивач, 7. еластична трака, 8. лимени држач казаљке, 9. казаљка, 10. назубљени лук, 11. лимени подешивач држача}

\textbf{Температурна компензација.} Померање биметалне траке зависи и од
температуре околине, што би утицало на тачност мерења. Зато је биметал у
облику слова U: један крак је радни, са загреваним влакном, а други је
компензациони и повезан је са постољем (слика 3.87). У почетном положају (I),
на обичној температури, оба крака су паралелна. Када порасте температура
околине, компензациони крак се савије удесно и помери тачку ослонца, а радни
крак се савије исто толико. Тако казаљка, без обзира на температуру околине,
остаје у истом положају.

\slika{0.45}{sl3-87.png}{Померање биметалне траке у облику слова U под утицајем температуре околине}

\textbf{Напонска компензација.} На тачност биметалних инструмената утиче и
стабилност напона напајања, па се уз њих обавезно уграђују мали биметални или
електронски стабилизатори напона.

\textbf{Показни инструмент са потенциометарским давачем.} Код давача са
потенциометром, као показни инструмент најчешће се користе инструменти са
укрштеним навојима и обртним магнетом, тзв. логометри. Немају опругу за
враћање: сваки навој ствара свој обртни момент, а покретни систем се
зауставља у положају у коме су ти моменти изједначени.

Ови инструменти имају три непомична навоја, од којих су два под правим углом,
а унутар њих је обртни магнет са казаљком. Шема везе дата је на слици 3.88.
Потенциометарски давач приказан је као променљиви отпорник (Ra).

У почетном положају, када је притисак нула, отпор давача је највећи, па кроз
трећи навој тече мала струја. Струја из акумулатора тече кроз први и други
навој, а магнет се поставља у правцу њиховог резултујућег поља, на почетак
скале. Када притисак расте, отпор давача се смањује, струја кроз трећи навој
расте, а са њом и поље $H_3$. Магнет се окреће ка новом резултујућем пољу
($H_{1,2} + H_3$), па је скретање сразмерно притиску. Отпорник $R_2$ служи за
температурну компензацију промене отпора навоја, а отпорником $R_1$ се
подешава јачина струје у навојима. Конструкција инструмента приказана је на
слици 3.89.

\slika{0.6}{sl3-88.png}{Шема везе инструмента са укрштеним навојима и обртним магнетом за мерење притиска са потенциометарским давачем}

\slika{0.5}{sl3-89.png}{Конструкција инструмента чија је шема везе дата на слици 3.88: 1. постоље, 2. осовина обртног магнета, 3. завртањ лежишта, 4 и 5. велики калем са 13 навојака, 6. заштитни омотач, 7. магнет који казаљку држи у нултом положају, 8. мали бочни калем, 9. крстасти држач, 10. завртањ, 11. матица лежишта, 12. матица, 13. навојак калема, 14. казаљка са контрателом}

%------------------------------------------------------------------
\subsubsection{Термометри за мерење температуре расхладне течности}
%------------------------------------------------------------------

За сигуран и економичан рад мотора, температура расхладне течности мора се
одржавати у одређеним границама, које зависе од конструкције и типа мотора.
Температуру аутоматски регулише термостат, а возач о њој добија информацију
преко мерног инструмента.

Температура се може мерити механички, на основу ширења неке течности, или
електрично, помоћу претварача који промену температуре претварају у
електрични сигнал. У даљем тексту обрађени су само електрични термометри.

\textbf{Принцип рада.} Слично је као код мерача притиска. Уз биметалне
инструменте користе се биметални давачи, а уз инструменте са укрштеним
навојима давачи чији се отпор мења са температуром. Код метала отпор расте са
температуром, а код полупроводника опада. Као давачи температуре најчешће се
користе полупроводници од керамичких једињења оксида тешких метала, познати
као \textbf{термистори}.

\textbf{Термометри са укрштеним навојем.} Најједноставнија варијанта има
два навоја под правим углом (слика 3.90). Давач (Ra) је термистор са
негативним температурним коефицијентом (NTC), везан паралелно са другим
навојем ($L_2$). Отпор термистора нагло опада са порастом температуре:

\begin{center}
\begin{tabular}{lccc}
    \toprule
    Температура [$^\circ$C] & 40  & 90   & 120 \\
    Отпорност [$\Omega$]     & 137 & 23,5 & 10,2 \\
    \bottomrule
\end{tabular}
\end{center}

Када је температура течности ниска, отпор давача је велики, па је струја у
паралелном навоју ($L_2$) скоро иста као у редно везаном навоју ($L_1$).
Магнет са казаљком поставља се у правцу резултујућег поља ($H_{rk}$), у коме
је поље другог навоја ($H_{2k}$) велико, а првог ($H_{1k}$) мало (слика 3.90в).
Са порастом температуре отпор давача опада, а са њим и напон на паралелном
делу везе (слика 3.90б). Струја у првом навоју ($L_1$) расте и поље $H_1$
се повећава, а струја у другом навоју ($L_2$) опада, па слаби и поље $H_2$.
Магнет са казаљком се окреће удесно, ка новом резултујућем пољу ($H_{rv}$).
Највеће скретање и ширина скале су релативно мали, што је главни недостатак
ових инструмената.

\slika{0.6}{sl3-90.png}{Шема везе инструмента са два укрштена навоја и термисторским давачем: а) електрична шема према стварном распореду навоја, предотпора, компензационог отпорника и отпорника давача, б) иста шема у развијеном облику, ради лакшег разумевања, в) векторски дијаграм}

На слици 3.91 дата је шема инструмента са три укрштена навоја, који је по
конструкцији сличан инструменту за мерење притиска (слика 3.88). Давач је
овде везан редно са трећим навојем (III). Отпорником ($R_e$) подешава се
напон напајања, а отпорником ($R_s$) скретање инструмента. Први и други навој
су под правим углом и кроз њих тече стална струја, која не зависи од отпора
давача. Резултујуће магнетно поље зато зависи само од струје у трећем
навоју, тј. од промене отпора давача. Мерно подручје инструмента је од
40 до 120\,$^\circ$C, а ширина скале $90^\circ$.

\slika{0.4}{sl3-91.png}{Шема везе инструмента са три укрштена навоја и термисторским давачем}

\textbf{Конструкција биметалног давача.} На слици 3.92 приказана је
конструкција биметалног давача за мерење температуре. Механизам давача
смештен је у месингану чауру (6), која се уврће у кућиште (1) шестоугаоним
завртњем. Биметална трака (3) једним крајем причвршћена је за изолациони
држач (7), а на другом крају носи помични контакт (4) и загревано влакно (2).
Један крај влакна спојен је са биметалом, а други, преко довода струје (8),
са прикључком (10). Непомични контакт (5) повезан је са масом преко чауре.

Када се укључи контакт-кључ, струја тече кроз инструмент, загревано влакно,
биметал и контакте до масе. У почетном положају, док је течност хладна,
контакти су затворени. Влакно загрева биметал, који се савија и отвара
контакте, затим се хлади и поново их затвара, па настају струјни импулси. Са
порастом температуре течности биметал додатно загрева и сама течност, па се
контакти брже отварају, а ефективна вредност струјних импулса опада.

\slika{0.5}{sl3-92.png}{Конструкција биметалног давача за мерење температуре: 1. кућиште, 2. загревано влакно, 3. биметална трака, 4. помични контакт, 5. непомични контакт, 6. месингана чаура, 7. изолациони држач, 8. довод струје, 9. изолатор, 10. електрични прикључак}

На слици 3.93 шематски је приказан рад биметалног давача са биметалним
инструментом. Инструмент ради у обрнутом смеру: при малим температурама
скретање је велико, а при великим мало. То је донекле и повољно, јер је током
вожње температура расхладне течности велика. У искљученом положају казаљка
стоји на 100\,$^\circ$C, тј. приближно на десном крају скале.

\slika{0.45}{sl3-93.png}{Шема везе биметалног давача са биметалним инструментом за мерење температуре расхладне течности}

\textbf{Термисторски давач.} На слици 3.94 приказана је конструкција
термисторског давача, који се најчешће везује са инструментом са укрштеним
навојима или са електронским инструментом. Плочица термистора (5) притиснута
је тањирастом опругом (4) и опружним проводником (2) на дно месинганог
кућишта (1), које је спојено са масом. Изолациона цев (3) одваја проводне
делове од кућишта. Отпор давача мења се, у зависности од температуре, од
око 1600 до 60~$\Omega$.

\slika{0.4}{sl3-94.png}{Попречни пресек термисторског давача за мерење температуре: 1. месингано кућиште, 2. опружни проводник струје, 3. изолациона цев, 4. тањираста опруга, 5. плоча термистора, 6. изолатор}

\textbf{Давач са отпорном жицом.} На слици 3.95 приказана је конструкција
температурног давача са отпорном металном жицом, који се користи код
камиона. Жица је намотана на тело (2) и смештена у месингану чауру (5).
Прикључна струја тече преко опруге (3) и отпорника ка месинганом кућишту и
маси. Електрични отпор давача при струји од 40~mA и температури од
80\,$^\circ$C износи $148{,}2 \pm 0{,}8~\Omega$.

\slika{0.5}{sl3-95.png}{Термички давач са отпорном металном жицом за камионе: 1. изолација, 2. метални отпорник, 3. проводна опруга, 4. заптивка, 5. месингано кућиште, 6. заштитна капица}

\textbf{Електронски термометар.} Поред класичних, користе се и
инструменти са електронским напајањем, укрштеним навојима и термисторским
давачем (слика 3.96). Када се укључи контакт-кључ (K), струја тече кроз
отпорник ($R_v$), потенциометар (P), термистор (Th) и отпорник ($R_c$). Пад
напона на потенциометру одређује струју базе транзистора (T), а од ње зависи
струја у колу инструмента. Пошто тај пад напона зависи од отпора термистора,
струја кроз инструмент директно зависи од температуре расхладне течности.

\slika{0.45}{sl3-96.png}{Шема везе електронског термометра}

%------------------------------------------------------------------
\subsubsection{Показивачи нивоа горива}
%------------------------------------------------------------------

Ови инструменти показују возачу стање горива у резервоару. Тачност им је
релативно мала, јер се од њих очекује само да покажу тренутну количину горива
у односу на запремину резервоара. Зато је подела на скали најчешће означена
бројевима 0, 1/4, 1/2, 3/4 и 1.

Код појединих возила на инструменту се пали мала сигнална лампица када ниво
горива падне на резерву. Возач тада зна приближно колико још километара
може да пређе.

\textbf{Биметални показивачи.} Код великог броја возила још се користе
биметални показивачи или инструменти са укрштеним навојима и
потенциометарским давачем. Биметални инструмент може имати
потенциометарски или биметални давач, а инструмент са укрштеним навојима
само потенциометарски.

Принцип рада биметалног показивача сличан је мерачу притиска (слика 3.85), с
тим што пловак у резервоару помера клизач по потенциометру, па је ефективна
вредност струјних импулса сразмерна нивоу горива. Највећи импулси су при пуном
резервоару, па је тада и скретање највеће. Принципијелна шема дата је на
слици 3.97. Иста је као код биметалних термометара, с тим што се исти
стабилизатор може користити за више мерних система.

\slika{0.45}{sl3-97.png}{Принципијелна шема везе биметалног показивача нивоа горива (VDO)}

\textbf{Показивачи са укрштеним навојима.} Код путничких возила ниже и средње
класе доста су распрострањени показивачи са укрштеним навојима и
потенциометарским давачем са пловком. Принципијелна шема дата је на слици
3.98.

Отпорници $R_1$ и $R_2$ су променљиви делови истог потенциометра: њихов збир
је увек исти, без обзира на положај клизача, а мења се само њихов однос. На
пример, $R_1$ је отпор од почетка потенциометра до клизача, а $R_2$ од клизача
до краја. Померањем клизача један отпорник се повећава, а други сразмерно
смањује. За разлику од термометра са слике 3.90, навоји су међусобно паралелно
везани, а сваки је редно везан са једним делом потенциометра.

Када је резервоар празан, пловак је доле, $R_1$ је најмањи, а $R_2$ највећи.
Струја у првом навоју ($L_1$) је највећа, а у другом ($L_2$) најмања. Поље
првог навоја ($H_1$) је много веће од поља другог ($H_2$), па се помично меко
гвожђе са казаљком поставља у правцу резултујућег поља, на почетак скале.
При пуном резервоару је обрнуто: пловак је горе, $R_1$ је највећи, а $R_2$
најмањи, па поље другог навоја преовлађује и меко гвожђе се поставља на крај
скале.

Иако је угао између навоја и њихових поља $90^\circ$, угао скретања ($\beta$)
између два резултујућа поља је око $45^\circ$, што даје релативно уску скалу.
Пошто се ниво горива ионако показује приближно, таква скала је сасвим
довољна. Казаљка се у искљученом положају враћа на нулу под дејством силе
теже, па јој је тада положај неодређен.

\slika{0.55}{sl3-98.png}{Принципијелна шема везе показивача нивоа горива са укрштеним навојима: а) електрична шема, б) развијена шема, в) дијаграм магнетних поља}

На слици 3.99 дата је конструкција и шематски приказ инструмента са укрштеним
навојима и помичним меким гвожђем. Струја се доводи преко прикључка (9),
тече кроз леви навој (1) и у тачки 12 се грана: кроз десни навој (7) и кроз
потенциометар давача (13) до масе. За разлику од претходног, овде су десни
навој и потенциометар везани паралелно. На слици 3.100 приказана је
конструкција потенциометарског давача са пловком.

\slika{0.6}{sl3-99.png}{Конструкција и шема везе инструмента са укрштеним навојима и помичним меким гвожђем: 1 и 7. навоји, 2. казаљка, 3. контактна плочица, 4, 6 и 11. држач меког гвожђа, 5. помично меко гвожђе, 8. држач, 9. изолатор, 10. електрични вод, 12. електрични прикључак, 13. потенциометар давача}

\slika{0.6}{sl3-100.png}{Конструкција потенциометарског давача са пловком у резервоару за гориво: 1. пловак, 2. осовина, 3 и 4. доњи и горњи део кућишта, 5. електрични прикључак, 6. спојне кукице, 7. полуга пловка, 8. спирална опруга за довод струје, 9. клизач, 10. изолационо тело, 11. електрични извод, 12. намотаји потенциометра}

\textbf{Електронски показивачи.} На слици 3.101 дата је шема електронског
показивача са потенциометарским давачем. Отпорници $R_1$ и $R_2$ чине
променљиви отпорник давача у мосној вези. Излазни напон давача појачава се
симетричним појачавачем са транзисторима $T_1$ и $T_2$, а мерни инструмент
везан је, као волтметар, између њихових колектора. Заједнички отпорник у колу
емитера обезбеђује стабилан рад. Колектор транзистора $T_2$ везан је за базу
транзистора $T_3$. Када ниво горива падне на минимум, напон између емитера и
базе транзистора $T_3$ достигне пробојни напон ценер-диоде ($Z_4$), $T_3$ се
отвара и пали црвену упозоравајућу лампицу на инструменту. Отпорником ($R_3$)
подешава се показивање инструмента при пуном резервоару, а отпорником ($R_4$)
при минималном нивоу (резерви), када се пали црвена контролна лампица.

\slika{0.55}{sl3-101.png}{Шема везе електронског показивача нивоа горива}

\textbf{Капацитивни давачи.} Код појединих возила, због облика резервоара,
не може се поставити потенциометарски или биметални давач са пловком, јер
захтева доста простора. Посебно место уградње тражи се и због безбедности у
случају судара. Тада се користе капацитивни или термички давачи, без
пловка.

Капацитивни давач користи промену капацитивности између облога. Диелектрична
константа горива је око два пута већа од ваздуха, па је и капацитивност
кондензатора чије су облоге у гориву око два пута већа него у ваздуху.
Капацитивност се зато мења сразмерно нивоу горива. Сонда давача прави се у
облику тракастог кондензатора, са облогама једна наспрам друге или једна поред
друге на еластичној фолији (слика 3.102а). Капацитивност се повећава већим
бројем облога, које се међусобно уплићу као чешаљ. Посебним обликовањем
трака (слика 3.102б) постиже се линеарно показивање и при неправилном облику
резервоара.

\slika{0.45}{sl3-102.png}{Капацитивни давачи: а) са равним облогама и б) са посебно обликованим облогама}

На слици 3.103 приказана је блок-шема електронског показивача нивоа горива са
капацитивним давачем. Кондензатори давача део су осцилатора, па одређују
његову учестаност ($f = 1/(K C_m) \approx 6$~kHz), која је сразмерна нивоу
горива. Претварач учестаност--напон претвара тај сигнал у напон, који покреће
мерни инструмент, а може да укључи и упозоравајућу лампицу за низак ниво
(резерву). Електронски део смешта се у кућиште давача на самом резервоару,
да би утицај капацитивности прикључних водова био што мањи.

\slika{0.5}{sl3-103.png}{Блок-шема електронског показивача нивоа горива са капацитивним давачем}

\textbf{Термички показивачи.} На слици 3.104 приказана је блок-шема
показивача са термичким давачима у резервоару. Термички давачи су
електрично загревани отпорници, чија температура зависи од тога колико су
уроњени у гориво, односно од нивоа горива. Пошто се њихов отпор мења са
температуром, при сталној струји пад напона на њима зависи од нивоа. Отпор
зависи и од температуре самог горива, што се компензује додатним отпорником у
резервоару. Он је од истог материјала, кроз њега тече иста струја, а
постављен је најниже, тако да је увек потопљен, па му отпор зависи само од
температуре горива. Разлика напона два отпорника се појачава и, преко
јединице за прилагођавање, доводи на инструмент.

Тачност оваквог мерења ремети таласање горива током вожње. Зато се мерна
сонда смешта у омотач у који гориво улази само кроз мали отвор на дну, па се
ниво око сонде не мења брзо.

\slika{0.5}{sl3-104.png}{Блок-шема електронског показивача нивоа горива са термичким давачем}

%------------------------------------------------------------------
\subsubsection{Амперметар и волтметар на возилу}
%------------------------------------------------------------------

Инсталација возила напаја се из акумулатора и алтернатора, везаних
паралелно. Током вожње алтернатор напаја све потрошаче, а ако је потрошња
мања, допуњава и акумулатор. Када је потрошња већа или мотор ради споро,
недостајућу енергију даје акумулатор, а када возило стоји, акумулатор је
једини извор. Без доброг стања акумулатора возило није у возном стању. Зато
се амперметром контролише струја пуњења и пражњења, а волтметром стање
акумулатора.

\paragraph{Амперметри, шема везе, опис и конструкција}

\textbf{Шема везе.} На возилима се најчешће уграђују амперметри са нултим
положајем на средини скале, тако да скретање казаљке улево или удесно
показује и смер струје. Обично је лево, са ознаком ($-$), пражњење, а десно,
са ознаком ($+$), пуњење акумулатора. Амперметар се везује између
алтернатора и акумулатора, према шеми на слици 3.105. Отпорник (R)
симболично представља све потрошаче, осим стартера, чија велика струја никад
не пролази кроз амперметар.

\slika{0.4}{sl3-105.png}{Шема везе амперметра на возилу}

\textbf{Опис и конструкција.} Због потребе за сталним померањем казаљке, на
возилима се најчешће уграђују амперметри са поларизованим меким гвожђем или
са обртним магнетом. На слици 3.106 приказана је конструкција амперметра са
поларизованим меким гвожђем и дијаграм магнетних поља.

Покретни систем чине казаљка (2) и меко гвожђе (6), на осовини (7) у
лежиштима (8). Меко гвожђе је од челика са малим губицима, па се лако
поларише и поставља у правцу резултујућег поља. У нултом положају држи га
стални магнет (4). Струја кроз инструмент ствара магнетно поље око масивног
проводника, чији смер зависи од смера струје. Поларизовано меко гвожђе
поставља се у правцу резултујућег поља сталног магнета и поља струје (слика
3.106б). При пражњењу акумулатора струја мења смер, па поље струје ($H_2$)
заузима супротан угао ($\varphi$) у односу на поље сталног магнета ($H_1$).
Меко гвожђе са казаљком тада скреће на другу страну, па амперметар показује
пражњење. Угао скретања зависи од јачине поља струје, тј. од саме струје.
Грешке због промене температуре компензују се магнетним отоком (3), слично
као код брзиномера.

\slika{0.4}{sl3-106.png}{Конструкција амперметра са поларизованим меким гвожђем и дијаграм магнетних поља: 1. електрични прикључак, 2. казаљка, 3. магнетни оток, 4. стални магнет, 5. постоље, 6. помично меко гвожђе, 7. осовина крећног система, 8. лежиште}

\textbf{Амперметар са обртним магнетом.} Конструкција је слична инструментима
са укрштеним навојима и обртним магнетом. Разлика је у томе што, уместо
малог навоја, стални магнет ствара поље управно на поље главног навоја. На
слици 3.107 приказан је принцип рада и шема везе.

Када кроз инструмент не тече струја, на обртни магнет (4) делује само
стални магнет (3), који га држи у нултом положају. Мањи, сразмерни део мерене
струје (већи део пролази кроз оток -- шант) пролази кроз навој (1) и ствара
магнетно поље, које се, као и код претходног амперметра, сабира са пољем
сталног магнета, а обртни магнет заузима положај у правцу резултујућег поља.
Помоћу отока мерни опсег се може знатно повећати, па се могу мерити и струје
до 100~A. Ради мањих губитака у водовима, због великих струја, оток се
поставља близу акумулатора, а до амперметра се води само мали, сразмерни
део струје.

\slika{0.4}{sl3-107.png}{Принцип рада амперметра са обртним магнетом и шема везе: 1. мерни навој, 2. струјни оток (шант), 3. стални магнет, 4. обртни магнет, 5. казаљка}

\paragraph{Волтметри, шема везе, опис и конструкција}

\textbf{Шема везе.} Волтметри се везују паралелно са акумулатором и
алтернатором. Када мотор не ради, показују стање напуњености акумулатора, а
када ради, напон пуњења (слика 3.109). Скала је обично подељена на жуто,
зелено и црвено поље, која посебно за искључен и укључен мотор означавају пун
или празан акумулатор, односно недовољно, добро и превелико пуњење (слика
3.108).

\slika{0.55}{sl3-108.png}{Показивање волтметра у зависности од стања и пуњења акумулатора}

\slika{0.4}{sl3-109.png}{Шема везе волтметра у електричној инсталацији возила}

\textbf{Опис и конструкција.} За волтметре се најчешће користе инструменти са
укрштеним навојима и обртним магнетом. На слици 3.110 приказана је
конструкција инструмента и дијаграм магнетних поља навоја. Обртни магнет је
на осовини са казаљком и увек се поставља у правцу резултујућег поља.
Разлика у односу на претходне инструменте је у томе што поље сталног магнета
у односу на поље навоја није под углом од $90^\circ$, већ од $135^\circ$, што
омогућава ширу скалу, од $90^\circ$.

У искљученом положају на обртни магнет делује само поље сталног магнета
($H_1$), па он стоји у нултом положају. У укљученом положају, при највећем
напону, кроз навој тече највећа струја, која ствара магнетно поље ($H_2$). Угао
између тих поља је $\varphi = 135^\circ$, а угао између резултујућег поља ($H_e$)
и нултог положаја је $\alpha = 90^\circ$, што је уједно и највеће скретање
инструмента.

\slika{0.5}{sl3-110.png}{Конструкција волтметра са укрштеним навојима и обртним магнетом: 1. стални магнет, 2. обртни магнет, 3. крупни намотаји, 4. доње постоље, 5. горње постоље, 6. осовина, 7. казаљка, $H_1$ -- поље сталног магнета, $H_2$ -- магнетно поље намотаја, $H_e$ -- резултујуће магнетно поље, $\varphi$ -- угао између поља сталног магнета и поља намотаја, $\alpha$ -- угао крајњег скретања}

%==================================================================
\subsection{Контролне лампице}
%==================================================================

Праћење и процена вредности појединих величина, битних за рад мотора и
погонску сигурност возила, не поверава се само пажњи возача. Тиме се смањује
његово оптерећење, па више пажње може да посвети условима саобраћаја, а
смањује се и потреба за стручним знањем за читање различитих инструмената.

Давачи који прате ове величине најчешће су конструисани као прекидачи. Они
мерене величине пореде са задатим референтним вредностима и, када их
достигну, укључују контролне лампице различитих боја, које упозоравају
возача. Најчешће се на овај начин прате притисак уља у мотору, температура
расхладне течности, пуњење акумулатора, ниво уља у кочионом систему и
слично. У посебну групу спадају контролне лампице које сигнализирају само
укључење неког уређаја.

\subsubsection{Контролна лампица за притисак уља у мотору}

Контролна лампица која сигнализира недозвољен пад притиска уља спада у
најстарије конструкције, а давачи за ту контролу су се до данас мало мењали.

\textbf{Давач.} На слици 3.111 дат је пресек давача контролне лампице за
притисак уља. Ради као обичан уљни прекидач, који укључује и искључује своје
контакте под утицајем притиска који делује на еластичну мембрану. Када је
притисак уља мањи од подешене вредности (обично 0,5~bar), мерна опруга (3),
преко завртња (2), потискује контактну плочицу (6) на кућиште (1), које је
спојено са масом, па контролна лампица светли. Када притисак уља порасте
изнад подешене вредности, мембрана (8) савлада опругу, раздвоји контакт
плочице (6) и масе, па се лампица гаси.

\slika{0.55}{sl3-111.png}{Давач контролне лампице за притисак уља: 1. метално кућиште са навојем, 2. завртањ за подешавање, 3. мерна опруга, 4. изолација, 5. контакт лампице, 6. контактна плочица, 7. тањираста опруга, 8. мембрана}

\subsubsection{Контролна лампица за температуру мотора}

Код мотора са унутрашњим сагоревањем једна од најважнијих величина која се
контролише је температура мотора. Најчешће се прати преко температуре
расхладне течности или уља, а код ваздушно хлађених мотора преко
температуре главе мотора. У новије време, код возила са Lambda-сондом и
катализатором, контролише се и температура издувних гасова. Давачи су
најчешће биметални контакти, а за веће температуре користе се давачи на
принципу промене отпора полупроводника или напона термоелемената.

\textbf{Биметални давачи.} На слици 3.112 дат је принцип рада биметалног давача
са контактом који затвара и са контактом који отвара при достигнутој
температури, а на слици 3.113 конструкција таквог давача.

\slika{0.5}{sl3-112.png}{Принцип рада биметалног давача са контактом који затвара и контактом који отвара при достигнутој температури}

Биметална трака (6) чврсто је спојена са месинганим телом (7) у врху
кућишта (2). Њен слободни крај, са помичним контактом (8), у искљученом
положају еластично је притиснут на непомични контакт (3), који је уједно и
завртањ за подешавање. Температура укључивања подешава се овим завртњем,
који се затим залије лаком (4), да се не би одврнуо. Струја се доводи преко
прикључка (9) и завртња до непомичног контакта, који је од кућишта изолован
изолационом чауром (1). Давачи са контактом који затвара производе се за
температуре од 92 до 100\,$^\circ$C, а са контактом који отвара за
температуре од 20\,$^\circ$C. (На слици је варијанта са контактом који отвара.)
Толеранција укључивања и искључивања је највише 5\,$^\circ$C.

\slika{0.55}{sl3-113.png}{Попречни пресек биметалног давача температуре (VDO): 1. изолациона чаура, 2. месингано кућиште, 3. завртањ за подешавање, 4. лак за заливање, 5. заштитна капа, 6. биметална трака, 7. месингана маса, 8. помоћни контакт, 9. електрични прикључак}

\subsubsection{Остале контролне лампице}

Контролне лампице могу се поделити у две групе. У прву групу спадају лампице
чије је паљење најважнија упозоравајућа информација за возача: њихово паљење
може значити да мотор не сме или не може да ради. То су већ описане лампице
за притисак уља, температуру мотора, пуњење акумулатора, ниво кочионе
течности, истрошеност кочионих облога и слично.

У другу групу спадају лампице које возача само подсећају да је неки уређај
укључен: контролне лампице за дуга и оборена светла, фарове за маглу,
упозоравајуће мигавце, грејач задњег стакла, ручну кочницу, чок (саух)
карбуратора итд. Оне се, за разлику од прве групе, пале паралелно са
укључењем уређаја. Код електричних уређаја (светла, грејачи) струја се доводи
директно од прекидача или релеја до лампице. Код механичких уређаја (ручна
кочница, чок) у механизам је уграђен мали микропрекидач, који пали лампицу
када се уређај активира.

Контролна лампица за мигавце објашњена је у поглављу о мигавцима, а
контролна лампица за пуњење акумулатора у поглављу о алтернатору.

\end{document}